%% ============================================================
%%  Section: Conclusion  (2026-04-07)
%%  Depositor Discipline in Modern Banking
%%  Target length: ~1 page (~350 words)
%%
%%  Improvements over conclusion_20260404.tex:
%%    - Opening paragraph no longer recites all six coefficients;
%%      states the core finding and what the design rules out.
%%    - Contribution paragraph added: positions against prior literature
%%      by naming the three identification problems the paper resolves.
%%    - Policy paragraph tightened; Correia et al. historical point
%%      retained but shortened.
%%    - Em dashes replaced throughout.
%%    - Ends with the broadest implication for the discipline-vs-stability tradeoff.
%%    - \label{sec:conclusion} resolves the undefined-reference warning in the intro.
%%
%%  Citations used (all in bib):
%%    correia2026bank, nier2006market, demirgucc2004market, flannery2019market
%% ============================================================

\section{Conclusion}
\label{sec:conclusion}

This paper provides causal evidence that uninsured depositors discipline community
banks in response to a credit-risk disclosure in normal times.  Exploiting the
mandatory CECL Day-One retained-earnings adjustment as a predetermined information
shock, we show that high-CECL community banks experienced significant contractions
in uninsured time deposits, simultaneous increases in insured deposits, higher
funding costs, and lower profitability relative to otherwise similar low-CECL
peers after 2023Q1 adoption.  A triple-difference specification with bank
$\times$ quarter fixed effects isolates a within-bank composition shift, ruling
out aggregate funding shocks as an alternative explanation.  The funding-cost and
profitability effects are amplified at banks in markets with financially
sophisticated depositor bases, consistent with informed depositors extracting
higher rates upon observing the CECL disclosure; the quantity effect shows a
gradual buildup in the event study but averages to an insignificant pooled
coefficient.

The findings resolve three problems that have contaminated prior tests of
depositor discipline.  Simultaneity---the co-evolution of bank risk and deposit
behavior in equilibrium---is addressed by the predetermined timing and magnitude
of the CECL shock.  Accounting opacity---the suppression of bank risk signals
under the incurred-loss standard---is the problem that CECL adoption itself
corrects.  Crisis-period confounding---the crowding-out of private discipline by
government guarantees in the samples with the most distress variation---is avoided
by focusing on a non-crisis, non-TBTF setting where theory predicts discipline
operates most freely.  Together, the design delivers estimates that are difficult
to attribute to anything other than informed uninsured depositors acting on newly
disclosed credit-risk information.

The policy relevance is immediate.  Congress is considering legislation that would
expand deposit insurance coverage for business accounts, most recently the Main
Street Depositor Protection Act introduced in March 2026.  Our estimates imply
that such an expansion would convert a substantial volume of currently sensitive
deposits into ones that, by design, carry no incentive to monitor.  The
disciplinary costs we document would not disappear; they would be transferred to
the Deposit Insurance Fund and borne across all insured institutions through
higher assessments.  This moral-hazard cost should be weighed against the
run-prevention benefits.  The historical record provides useful perspective on
those benefits: \citet{correia2026bank}, examining over 3{,}400 bank runs in
the United States between 1863 and 1934, find that failures following runs were
concentrated almost entirely at institutions with poor fundamentals, while
well-capitalized banks survived runs through liquidity support.  If runs are
disciplinary rather than purely destabilizing, then policies that eliminate the
threat of runs at weak institutions remove a governance mechanism from the banks
that need it most.  The right policy framework should preserve the disciplinary
pressure that uninsured depositors provide, while extending stability protections
in ways that do not wholesale eliminate the incentive to monitor.
